\chapter{Criterio de comparabilidad}

\section{Comparabilidad disciplinada}
\begin{definition}
Dos formulaciones, configuraciones o trayectorias son comparables en sentido disciplinado cuando pueden ponerse en relación dentro de un mismo dominio declarado sin pérdida inmediata de herencia canónica, legibilidad estructural o nivel de tratamiento.
\end{definition}

\begin{observation}
La comparabilidad disciplinada no equivale a identidad ni a equivalencia plena. Solo fija la posibilidad de relación ordenada bajo criterios compartidos.
\end{observation}

\section{Condición mínima de comparabilidad}
\begin{definition}
Existe comparabilidad mínima cuando concurren simultáneamente:
\begin{enumerate}
  \item dependencia respecto de un mismo suelo canónico o de suelos compatibles;
  \item compatibilidad de nivel jerárquico suficiente;
  \item umbral operativo no contradictorio;
  \item posibilidad de lectura dentro de una misma ventana declarada o de ventanas traducibles entre sí.
\end{enumerate}
\end{definition}

\begin{observation}
Sin estas condiciones, la comparación corre el riesgo de confundir diferencia estructural con mera heterogeneidad de marco.
\end{observation}

\section{Comparabilidad y transición}
\begin{definition}
Si dos formulaciones pertenecen a niveles distintos, solo son comparables de manera disciplinada cuando la transición entre niveles puede declararse sin ruptura de herencia ni pérdida ilegible de estructura.
\end{definition}

\begin{observation}
La comparabilidad entre niveles no exige igualdad de fuerza, pero sí una relación trazable entre los términos comparados.
\end{observation}

\section{Comparabilidad y proyección}
\begin{definition}
Dos piezas pueden ser comparadas mediante lectura proyectada solo si la proyección conserva la información mínima requerida para el propósito comparativo declarado.
\end{definition}

\begin{observation}
La proyección puede facilitar comparación, pero no autoriza a ignorar información estructural cuya pérdida vuelva equívoca la relación.
\end{observation}

\section{Condición mínima de no comparabilidad}
\begin{definition}
Se considera que dos piezas no son comparables en sentido disciplinado cuando:
\begin{enumerate}
  \item pertenecen a marcos incompatibles;
  \item no comparten dependencia canónica suficiente;
  \item operan en niveles no traducibles;
  \item la relación entre ellas solo puede formularse mediante retórica o analogía sin soporte estructural.
\end{enumerate}
\end{definition}

\begin{observation}
Esta cláusula protege el canon intermedio frente a comparaciones infladas o artificiales.
\end{observation}

\section{Cláusula de alcance}
\begin{observation}
El criterio de comparabilidad fija una regla mínima de relación entre piezas del corpus. No constituye todavía una teoría completa de equivalencia, orden total o clasificación exhaustiva de semejanzas y diferencias.
\end{observation}

